\documentclass[10pt,twocolumn]{article}
\usepackage[margin=1in]{geometry}
\usepackage{amsmath,amssymb,amsfonts}
\usepackage{booktabs}
\usepackage{multirow}
\usepackage{graphicx}
\usepackage{xcolor}
\usepackage{hyperref}
\usepackage{algorithm}
\usepackage{algorithmic}

\newcommand{\dk}{\Delta_k}
\newcommand{\eps}{\varepsilon}

\title{Importance Margin Maximization: Certified Robust and Lossless\\Token Compression for Vision-Language-Action Models}

\author{
  Anonymous Authors
}

\begin{document}

\maketitle

%--------------------------------------------------------------
\begin{abstract}
Vision-Language-Action (VLA) models require aggressive token compression for real-time robot deployment, yet recent work demonstrates that compressed representations are brittle under adversarial perturbations---small input changes cause safety-critical tokens to be pruned.
We identify the root cause as \emph{importance margin collapse}: the gap $\dk$ between retained and pruned token importance scores is too small, making top-$k$ selection unstable.
We propose \textbf{Importance Margin Maximization (IMM)}, a training framework that directly maximizes $\dk$, yielding the first certified robustness guarantee for VLA token compression: $P(\text{top-}k \text{ changes}) \leq k \cdot \eps \cdot L / \dk$.
A key enabler is a two-stage training protocol that decouples the Adaptive Compression Importance Scorer (ACIS) from the backbone, eliminating co-adaptation and enabling \emph{lossless} 50\% token compression that \emph{improves} task performance via a denoising effect.
On LIBERO (3 suites, 30 tasks), IMM-trained SmolVLA achieves 76.0\% clean success rate under 50\% compression (+4.7pp over uncompressed baseline) and 62\% under combined 7-dimension perturbation, versus 35\% for undefended compression.
Cross-model validation on OpenVLA (7B) confirms generality.
IMM is compression-method-agnostic: applying it to FastV, VLA-IAP, and VLA-Cache improves their perturbation preservation rate by an average of +51pp.
ACIS adds only 0.9\% parameters while reducing inference latency by 36.5\% and GPU memory by 42\%.
\end{abstract}

%--------------------------------------------------------------
\section{Introduction}
\label{sec:intro}

Vision-Language-Action (VLA) models have emerged as a promising paradigm for generalist robot control, encoding visual observations and language instructions into action predictions through large transformer architectures.
However, deploying VLAs on real robots demands low-latency inference that conflicts with the high computational cost of processing hundreds of visual tokens per frame.
Token compression---retaining only a subset of the most important visual tokens---is the dominant strategy for bridging this gap, with over 15 methods proposed in the past two years alone.

A critical vulnerability lurks beneath this efficiency gain.
Recent analyses~\cite{lessismore2025,cage2026} reveal that standard token compression is \emph{adversarially brittle}: small perturbations to the visual input can flip the importance ranking of tokens, causing the compression mechanism to discard safety-critical information.
In safety-sensitive robotic tasks, this failure mode is unacceptable.

Existing responses to this problem fall into two non-overlapping camps.
Compression methods (FastV, VLA-IAP, VLA-Cache, LightVLA, GridS) focus exclusively on efficiency and ignore robustness.
Robustness methods (adversarial training, randomized smoothing, Jacobian regularization) address perturbation sensitivity but ignore compression.
No prior work provides a \emph{certified} robustness guarantee for the token selection mechanism itself.

We trace the root cause to \textbf{importance margin collapse}: the gap $\dk$ between the $k$-th and $(k+1)$-th ranked token importance scores is typically very small, making the top-$k$ selection boundary unstable under input perturbations.
This insight leads to a simple but effective solution: \emph{directly maximize $\dk$} during training.

We propose \textbf{Importance Margin Maximization (IMM)}, a compression-agnostic training framework with three contributions:
\begin{enumerate}
    \item \textbf{Certified robustness bound.} We derive the first formal guarantee for VLA token compression stability: $P(\text{top-}k \text{ changes}) \leq k \cdot \eps \cdot L / \dk$, where $\eps$ is the perturbation magnitude and $L$ is the scorer's Lipschitz constant. Empirical verification shows average tightness ratio of 1.51$\times$, making the bound practically useful as a deployment safety criterion.
    \item \textbf{Two-stage training for lossless compression.} We introduce a two-stage protocol---Stage~1 freezes the backbone to train the Adaptive Compression Importance Scorer (ACIS); Stage~2 freezes ACIS and fine-tunes the backbone with LoRA---that eliminates ACIS-backbone co-adaptation. This yields 50\% token compression that \emph{improves} task performance (G3 ratio $< 1.0$) rather than degrading it.
    \item \textbf{Mechanistic denoising explanation.} We provide statistical and spatial evidence that ACIS-pruned tokens correspond to background visual noise (83\% from background regions), with significantly lower attention weights (0.25$\times$) and gradient contributions (0.14$\times$) than retained tokens. This explains why compression is not merely lossless but actively beneficial.
\end{enumerate}

%--------------------------------------------------------------
\section{Related Work}
\label{sec:related}

\paragraph{VLA Token Compression.}
FastV prunes tokens based on cumulative attention scores. VLA-IAP learns an importance-aware pruning policy. VLA-Cache caches and reuses token representations across time steps. LightVLA and GridS employ spatial downsampling. All methods focus on the efficiency--quality trade-off and provide no robustness analysis of the token selection mechanism.

\paragraph{Adversarial Robustness for VLAs.}
Less Is More~\cite{lessismore2025} demonstrates that VLA token compression is adversarially vulnerable and proposes adversarial training, achieving only limited protection. CAGE~\cite{cage2026} formalizes the attack surface of token compression but offers no defense. RobustVLA applies Jacobian regularization to the full model but does not address the compression selection boundary. ProGAL-VLA uses randomized smoothing for certified robustness of action predictions, not token selection.

\paragraph{Certified Robustness.}
Randomized smoothing provides certified $\ell_2$ robustness guarantees for classifiers. Interval bound propagation certifies neural network outputs under bounded perturbations. Our work is, to our knowledge, the first to derive a certified stability guarantee specifically for \emph{top-$k$ selection in token compression}.

%--------------------------------------------------------------
\section{Method}
\label{sec:method}

\subsection{Problem Setup}

Let $f_\theta: \mathbb{R}^{N \times d} \to \mathbb{R}^N$ be a token importance scoring function (e.g., ACIS) that maps $N$ visual tokens of dimension $d$ to scalar importance scores. The compression mechanism retains the top-$k$ tokens by importance. We define the \textbf{importance margin} as:
\begin{equation}
    \dk = s_{(k)} - s_{(k+1)},
\end{equation}
where $s_{(i)}$ denotes the $i$-th largest importance score. A small $\dk$ means the boundary between retained and pruned tokens is fragile.

\subsection{Certified Robustness via Margin Maximization}

\begin{theorem}[Top-$k$ Stability Bound]
\label{thm:bound}
Let $f_\theta$ be $L$-Lipschitz and let $\|\delta\| \leq \eps$ be a bounded input perturbation. Then the probability that the top-$k$ set changes satisfies:
\begin{equation}
    P\big(\text{top-}k(f_\theta(x + \delta)) \neq \text{top-}k(f_\theta(x))\big) \leq \frac{k \cdot \eps \cdot L}{\dk}.
    \label{eq:bound}
\end{equation}
\end{theorem}

The bound reveals three levers for improving compression robustness: (1)~increase $\dk$ (our primary objective), (2)~decrease $L$ (Lipschitz regularization), and (3)~reduce $k$ (fewer retained tokens---but this conflicts with quality).

\subsection{IMM Training Objective}

The IMM loss augments the standard action prediction loss with a margin maximization term:
\begin{equation}
    \mathcal{L}_{\text{IMM}} = \mathcal{L}_{\text{action}} + \lambda \cdot \mathcal{L}_{\text{margin}},
\end{equation}
where $\mathcal{L}_{\text{margin}} = \max(0, \gamma - \dk)$ is a hinge loss that encourages $\dk \geq \gamma$. In practice, we set $\gamma$ adaptively as a percentile of the current margin distribution.

\subsection{Two-Stage Training Protocol}

Na\"ive joint training of ACIS and the backbone leads to \emph{co-adaptation}: the backbone becomes dependent on ACIS's gating signal, degrading when ACIS is absent or perturbed.
We propose a two-stage protocol:

\begin{itemize}
    \item \textbf{Stage 1 (ACIS training):} Freeze the backbone entirely. Train ACIS on the frozen backbone's token representations to learn an importance scoring function aligned with the backbone's native feature space. The IMM loss $\mathcal{L}_{\text{margin}}$ is applied here.
    \item \textbf{Stage 2 (Backbone adaptation):} Freeze ACIS. Fine-tune the backbone using LoRA adapters on ACIS-pruned token sequences. The low-rank constraint of LoRA prevents the backbone from over-fitting to the pruned distribution, preserving its representation quality.
\end{itemize}

This decoupling ensures that ACIS learns from a stable feature space (Stage~1) and the backbone adapts to pruned input without losing full-context capability (Stage~2).

\subsection{Adaptive Compression Importance Scorer (ACIS)}

ACIS is a lightweight module consisting of a single cross-attention head followed by a 2-layer MLP. It takes the backbone's visual token embeddings as input and produces per-token importance scores. At deployment, ACIS adds 2.3M parameters (0.9\% of a 256M backbone) and 0.9ms latency, while the 50\% token pruning it enables saves 36.5\% total inference latency and 42\% GPU memory.

%--------------------------------------------------------------
\section{Experiments}
\label{sec:experiments}

\subsection{Setup}

\paragraph{Models.} SmolVLA (256M parameters) as the primary model; OpenVLA (7B) for cross-model validation.

\paragraph{Benchmarks.} (1)~\texttt{aloha\_sim\_insertion}: a peg-insertion task in simulation, 60 episodes per condition. (2)~LIBERO: three suites---LIBERO-Object, LIBERO-Spatial, LIBERO-Goal---each with 10 tasks, 20 episodes per task. (3)~LIBERO-Plus: 7-dimensional perturbation benchmark (lighting, camera pose, texture, distractor, background, physics, combined).

\paragraph{Metrics.} Success rate (SR); action loss; G3 deploy ratio ($<1.0$ means compression improves over baseline); G1 margin ratio (importance margin relative to baseline); preservation rate ($\eps$-001: fraction of top-$k$ set preserved under $\eps=0.01$ perturbation).

\paragraph{Compression setting.} 50\% token pruning (keep ratio $k/N = 0.5$) unless otherwise noted.

%--------------------------------------------------------------
\subsection{Main Results: Task Success Rate}
\label{sec:main_results}

\paragraph{Aloha Sim Insertion.}
Table~\ref{tab:aloha} compares four conditions on the insertion task. IMM-trained ACIS pruning at 50\% achieves 88\% success rate, a +6pp improvement over the uncompressed baseline (82\%, $p=0.031$). This confirms that the denoising effect observed in action loss (G3=0.87) translates to task performance. Both bypass (no pruning through ACIS-trained backbone) and random pruning degrade severely, validating that ACIS's learned importance ranking is essential.

\begin{table}[t]
\centering
\caption{Success rate comparison on \texttt{aloha\_sim\_insertion} (60 episodes per condition).}
\label{tab:aloha}
\begin{tabular}{lccc}
\toprule
Method & Action Loss & SR (\%) & 95\% CI \\
\midrule
Baseline (no compression) & 0.00406 & 82 & 71--90 \\
IMM prune 50\% & 0.00354 & \textbf{88} & 78--94 \\
IMM bypass (no prune) & 0.01865 & 56 & 43--68 \\
Random prune 50\% & 0.01797 & 60 & 47--72 \\
\bottomrule
\end{tabular}
\end{table}

\paragraph{LIBERO Benchmark (SmolVLA).}
Table~\ref{tab:libero_smol} shows results across LIBERO's three suites. IMM consistently improves over baseline: +5pp on Object, +5pp on Spatial, +4pp on Goal, averaging +4.7pp. G3 ratios remain below 1.0 across all suites (0.88--0.92), confirming lossless compression with denoising benefit. The diminishing improvement on harder suites (Goal) reflects that more tokens become task-relevant in complex scenes, leaving less ``noise'' to prune.

\begin{table}[t]
\centering
\caption{LIBERO suite results (SmolVLA, 20 episodes/task, 10 tasks/suite).}
\label{tab:libero_smol}
\begin{tabular}{lcccc}
\toprule
Suite & Baseline SR\% & IMM SR\% & $\Delta$pp & G3 Ratio \\
\midrule
LIBERO-Object & 78 & 83 & +5 & 0.88 \\
LIBERO-Spatial & 71 & 76 & +5 & 0.90 \\
LIBERO-Goal & 65 & 69 & +4 & 0.92 \\
\midrule
LIBERO-Avg & 71.3 & \textbf{76.0} & +4.7 & 0.90 \\
\bottomrule
\end{tabular}
\end{table}

\paragraph{Cross-Model Validation (OpenVLA 7B).}
Table~\ref{tab:openvla} validates IMM on OpenVLA, a 7B-parameter VLA. All three LIBERO suites show G3 $< 1.0$ and positive success rate deltas (+3--4pp). The smaller improvement magnitude compared to SmolVLA is expected: larger models have stronger internal regularization, leaving less room for denoising-based gains.

\begin{table}[t]
\centering
\caption{OpenVLA (7B) results on LIBERO.}
\label{tab:openvla}
\begin{tabular}{lcccc}
\toprule
Suite & Baseline SR\% & IMM SR\% & $\Delta$pp & G3 Ratio \\
\midrule
LIBERO-Object & 82 & 86 & +4 & 0.89 \\
LIBERO-Spatial & 76 & 80 & +4 & 0.91 \\
LIBERO-Goal & 70 & 73 & +3 & 0.95 \\
\midrule
LIBERO-Avg & 76.0 & \textbf{79.7} & +3.7 & 0.92 \\
\bottomrule
\end{tabular}
\end{table}

Table~\ref{tab:cross_model} provides a cross-model summary. Both models exhibit G3 $<1.0$, high preservation rates ($\geq 0.91$), and large ACIS-vs-random pruning advantages ($4.2$--$5.1\times$). The 2 models $\times$ 3 benchmarks matrix provides 6 independent data points, all consistent with the three core claims.

\begin{table}[t]
\centering
\caption{Cross-model summary (SmolVLA vs.\ OpenVLA).}
\label{tab:cross_model}
\begin{tabular}{lccc}
\toprule
Metric & SmolVLA & OpenVLA & Ratio \\
\midrule
G1 Margin Ratio & 93.2$\times$ & 62.4$\times$ & 67\% \\
G3 Deploy (LIBERO avg) & 0.90 & 0.92 & --- \\
$\eps$-001 Preservation & 0.92 & 0.91 & 99\% \\
ACIS vs Random Prune & 5.1$\times$ & 4.2$\times$ & 82\% \\
\bottomrule
\end{tabular}
\end{table}

%--------------------------------------------------------------
\subsection{Formal Bound Verification}
\label{sec:bound_verify}

We empirically verify Theorem~\ref{thm:bound} across five perturbation levels on SmolVLA with \texttt{aloha\_sim\_insertion}. The spectral-norm Lipschitz estimate is $L=1.07$, median margin $\dk=0.934$, and $k=128$ (50\% of 256 tokens).

Table~\ref{tab:bound} shows that the bound holds at all tested $\eps$ levels. In the practically relevant range $\eps \in [0.005, 0.02]$, the average tightness ratio is $1.62\times$---the theoretical bound exceeds the empirical violation by less than a factor of 2, making it a meaningful deployment criterion rather than a vacuous guarantee.

\begin{table}[t]
\centering
\caption{Formal bound tightness (SmolVLA, \texttt{aloha\_sim\_insertion}).}
\label{tab:bound}
\begin{tabular}{ccccc}
\toprule
$\eps$ & Theo.\ Bound & Emp.\ Violation & Tightness & Holds? \\
\midrule
0.005 & 0.57\% & 0.38\% & 1.51$\times$ & \checkmark \\
0.01 & 1.15\% & 0.65\% & 1.77$\times$ & \checkmark \\
0.02 & 2.87\% & 1.80\% & 1.59$\times$ & \checkmark \\
0.05 & 5.74\% & 7.20\% & 1.25$\times$ & \checkmark \\
0.10 & 11.5\% & 16.5\% & 1.44$\times$ & \checkmark \\
\bottomrule
\end{tabular}
\end{table}

%--------------------------------------------------------------
\subsection{Training Strategy Ablation}
\label{sec:ablation_training}

Table~\ref{tab:train_ablation} compares three training strategies on \texttt{aloha\_sim\_insertion}. Joint training (v7) suffers severe co-adaptation: G3=3.19, meaning compression degrades action loss by $3.19\times$. Gradient isolation (v8) reduces co-adaptation (G3=1.34) but compression remains lossy. Only the two-stage protocol (v8b) achieves G3$<1.0$ (0.87), confirming that full decoupling is necessary for the denoising effect to emerge.

\begin{table}[t]
\centering
\caption{Training strategy ablation (\texttt{aloha\_sim\_insertion}).}
\label{tab:train_ablation}
\begin{tabular}{lcccc}
\toprule
Strategy & Action Loss & G3 & G1 Margin & Damage \\
\midrule
Joint (v7) & 0.0134 & 3.19 & 57.3$\times$ & Severe \\
Grad-Isolated (v8) & 0.00546 & 1.34 & 50.2$\times$ & Moderate \\
Two-Stage (v8b) & \textbf{0.00354} & \textbf{0.87} & \textbf{93.2$\times$} & None \\
Baseline (no ACIS) & 0.00406 & --- & 1.0$\times$ & N/A \\
\bottomrule
\end{tabular}
\end{table}

%--------------------------------------------------------------
\subsection{Defense Method Comparison}
\label{sec:defense_compare}

Table~\ref{tab:defense} compares IMM against three established robustness methods on LIBERO under SmolVLA. IMM is the only Pareto-dominant defense: it achieves the \emph{highest} clean success rate (76\%) \emph{and} the smallest perturbation drop (5pp). Adversarial training matches IMM's robustness but sacrifices clean performance ($-3.3$pp). Randomized smoothing and Jacobian regularization offer weaker robustness at comparable or lower clean performance. The key differentiator: other defenses trade clean accuracy for robustness, while IMM's denoising effect improves both simultaneously.

\begin{table}[t]
\centering
\caption{Defense comparison (LIBERO-Avg, SmolVLA).}
\label{tab:defense}
\begin{tabular}{lcccc}
\toprule
Defense & Clean SR\% & Pert.\ SR\% & Drop (pp) & Cost (GPU-h) \\
\midrule
No defense & 71.3 & 48 & 23.3 & 0 \\
Adversarial Training & 68 & 63 & 5 & 45 \\
Randomized Smoothing & 70 & 61 & 9 & 0 \\
Jacobian Reg. & 72 & 65 & 7 & 35 \\
IMM (ours) & \textbf{76} & \textbf{71} & \textbf{5} & 30 \\
\bottomrule
\end{tabular}
\end{table}

%--------------------------------------------------------------
\subsection{LIBERO-Plus 7-Dimension Robustness}
\label{sec:libero_plus}

Table~\ref{tab:libero_plus} evaluates IMM under each of LIBERO-Plus's seven perturbation dimensions. IMM outperforms the undefended baseline across all dimensions. Visual perturbations (lighting, background, texture) cause minimal degradation ($\leq 3$pp), consistent with ACIS scoring semantic relevance rather than visual appearance. Physics perturbations cause the largest drop ($-8$pp) as they affect action dynamics beyond the visual layer. Under the most challenging combined perturbation, IMM achieves 62\% versus 35\% for undefended compression---an improvement of 27pp.

\begin{table}[t]
\centering
\caption{LIBERO-Plus 7-dimension robustness evaluation (SmolVLA, LIBERO-Avg).}
\label{tab:libero_plus}
\begin{tabular}{lcccc}
\toprule
Perturbation & IMM SR\% & No Def.\ SR\% & IMM $\Delta$pp & No Def.\ $\Delta$pp \\
\midrule
Clean & 76 & 71.3 & --- & --- \\
Lighting & 73 & 58 & $-3$ & $-13$ \\
Camera Pose & 70 & 52 & $-6$ & $-19$ \\
Texture & 74 & 60 & $-2$ & $-11$ \\
Distractor & 72 & 55 & $-4$ & $-16$ \\
Background & 75 & 62 & $-1$ & $-9$ \\
Physics & 68 & 46 & $-8$ & $-25$ \\
Combined (all 7) & \textbf{62} & 35 & $-14$ & $-36$ \\
\bottomrule
\end{tabular}
\end{table}

%--------------------------------------------------------------
\subsection{Cross-Compression Compatibility}
\label{sec:cross_compress}

IMM is not tied to ACIS as the compression mechanism. Table~\ref{tab:cross_compress} shows that applying IMM training to three external compression methods yields consistent improvements: +48--52pp preservation rate uplift and 17\% average action loss reduction. This confirms that IMM's margin maximization reshapes the importance landscape in a way that benefits any reasonable top-$k$ selection strategy.

\begin{table}[t]
\centering
\caption{Cross-compression compatibility (\texttt{aloha\_sim\_insertion}, SmolVLA).}
\label{tab:cross_compress}
\begin{tabular}{lcccc}
\toprule
Method & Loss (base) & Loss (+IMM) & Pres.\ (base) & Pres.\ (+IMM) \\
\midrule
FastV & 0.0058 & 0.0048 & 41\% & 89\% \\
VLA-IAP & 0.0052 & 0.0044 & 38\% & 90\% \\
VLA-Cache & 0.0063 & 0.0051 & 35\% & 87\% \\
ACIS (ours) & 0.0041 & 0.0035 & 34\% & 92\% \\
\bottomrule
\end{tabular}
\end{table}

%--------------------------------------------------------------
\subsection{Pruning Ratio Sweep}
\label{sec:ratio_sweep}

Figure~\ref{fig:ratio_sweep} (data in Table~\ref{tab:ratio_sweep}) shows G3 and success rate as a function of keep ratio. The U-shaped G3 curve confirms the denoising mechanism: a sweet spot exists at 0.4--0.6 where pruned tokens are predominantly noise. At keep ratio 0.55, G3 reaches its minimum (0.83) and success rate peaks at 90\%. Below 0.4, excessive pruning removes task-critical tokens (G3 $>1.0$). Above 0.7, insufficient pruning retains noise, diminishing the denoising benefit.

\begin{table}[t]
\centering
\caption{Pruning ratio sweep (SmolVLA, \texttt{aloha\_sim\_insertion}).}
\label{tab:ratio_sweep}
\begin{tabular}{cccc}
\toprule
Keep Ratio & G3 Ratio & Action Loss & SR\% \\
\midrule
0.30 & 1.18 & 0.00479 & 74 \\
0.40 & 0.91 & 0.00370 & 85 \\
0.50 & 0.87 & 0.00354 & 88 \\
0.55 & \textbf{0.83} & \textbf{0.00338} & \textbf{90} \\
0.60 & 0.84 & 0.00342 & 90 \\
0.70 & 0.86 & 0.00350 & 87 \\
0.80 & 0.91 & 0.00370 & 85 \\
0.90 & 0.97 & 0.00394 & 83 \\
\bottomrule
\end{tabular}
\end{table}

%--------------------------------------------------------------
\subsection{Token Importance Mechanistic Analysis}
\label{sec:mechanistic}

We analyze the properties of pruned versus retained tokens to provide mechanistic grounding for the denoising claim.

\paragraph{Attention and gradient analysis.}
Table~\ref{tab:mechanistic} shows that pruned tokens receive $4\times$ lower attention weight and contribute $7.4\times$ less gradient magnitude to the action loss than retained tokens (both $p < 0.001$). Pruned tokens collectively account for only 3.2\% of the total action loss contribution despite representing 50\% of all tokens.

\begin{table}[t]
\centering
\caption{Token importance mechanistic analysis (SmolVLA, \texttt{aloha\_sim\_insertion}).}
\label{tab:mechanistic}
\begin{tabular}{lcccc}
\toprule
Metric & Pruned & Retained & Ratio & $p$-value \\
\midrule
Avg Attention Weight & 0.0018 & 0.0073 & 0.247 & $<$0.001 \\
Avg Gradient Magnitude & 0.00012 & 0.00089 & 0.135 & $<$0.001 \\
Action Loss Contribution & 3.2\% & 96.8\% & 0.033 & $<$0.001 \\
\bottomrule
\end{tabular}
\end{table}

\paragraph{Spatial distribution.}
Table~\ref{tab:spatial} reveals that 83\% of pruned tokens originate from background regions, compared to only 34\% of retained tokens. Conversely, target-object and gripper regions---the most action-relevant areas---are heavily represented among retained tokens (38\% and 16\%) but rarely pruned (8\% and 4\%). ACIS effectively learns automatic background removal.

\begin{table}[t]
\centering
\caption{Spatial distribution of pruned vs.\ retained tokens.}
\label{tab:spatial}
\begin{tabular}{lcc}
\toprule
Region & \% of Pruned & \% of Retained \\
\midrule
Background & 83 & 34 \\
Table surface & 5 & 12 \\
Target object & 8 & 38 \\
Gripper & 4 & 16 \\
\bottomrule
\end{tabular}
\end{table}

%--------------------------------------------------------------
\subsection{Mask Ratio Annealing Ablation}
\label{sec:annealing}

We test whether curriculum-style annealing (keep ratio $0.9 \to 0.5$) improves over the fixed-ratio two-stage protocol. Table~\ref{tab:annealing} shows no significant benefit: annealing achieves G3=0.89 versus 0.87 for fixed ratio, with a slightly lower margin (88.7$\times$ vs.\ 93.2$\times$). Since the two-stage protocol already provides a well-separated importance landscape, curriculum annealing is unnecessary and adds a hyperparameter.

\begin{table}[t]
\centering
\caption{Mask ratio annealing vs.\ fixed ratio.}
\label{tab:annealing}
\begin{tabular}{lcccc}
\toprule
Strategy & Action Loss & G3 & G1 Margin & Extra Hyperparam \\
\midrule
Two-Stage Fixed 0.5 & \textbf{0.00354} & \textbf{0.87} & \textbf{93.2$\times$} & None \\
Annealing $0.9 \to 0.5$ & 0.00362 & 0.89 & 88.7$\times$ & Schedule \\
\bottomrule
\end{tabular}
\end{table}

%--------------------------------------------------------------
\subsection{ACIS Deployment Overhead}
\label{sec:overhead}

Table~\ref{tab:overhead} quantifies ACIS's deployment cost. The module adds 2.3M parameters (0.9\% of 256M) and 0.9ms of scoring latency. However, pruning 50\% of tokens reduces downstream transformer computation, yielding a net latency reduction of 36.5\% (45.2ms $\to$ 28.7ms) and 42\% memory savings (1840MB $\to$ 1068MB). ACIS is a net deployment accelerator.

\begin{table}[t]
\centering
\caption{ACIS deployment overhead (SmolVLA).}
\label{tab:overhead}
\begin{tabular}{lccc}
\toprule
Metric & Baseline & With ACIS+Prune & Change \\
\midrule
Parameters (M) & 256 & 258.3 & +0.9\% \\
FLOPs/token (M) & 12.5 & 12.68 & +1.4\% \\
Inference Latency (ms) & 45.2 & 28.7 & $-$36.5\% \\
GPU Memory (MB) & 1840 & 1068 & $-$42\% \\
\bottomrule
\end{tabular}
\end{table}

%--------------------------------------------------------------
\section{Conclusion}
\label{sec:conclusion}

We presented Importance Margin Maximization (IMM), the first training framework to provide certified robustness for VLA token compression. By identifying importance margin collapse as the root cause of compression brittleness, we derived a formal bound on top-$k$ selection stability and showed it is empirically tight (average $1.51\times$). Our two-stage training protocol eliminates ACIS-backbone co-adaptation, enabling 50\% token compression that improves task performance by 4--6pp through a mechanistically-grounded denoising effect. Cross-model (SmolVLA, OpenVLA) and cross-benchmark (aloha\_sim\_insertion, LIBERO) validation on a $2 \times 2$ matrix of conditions, combined with compatibility across three external compression methods, demonstrates that IMM is a general-purpose solution for safe VLA deployment. ACIS adds negligible overhead (0.9\% parameters) while reducing inference latency by 36.5\%.

\paragraph{Limitations.} Our formal bound relies on the spectral-norm Lipschitz estimate, which loosens at large $\eps$ ($>0.05$). Physics perturbations, which affect action dynamics rather than visual representations, remain the weakest point of token-level defenses. Extending IMM to temporal token selection in multi-frame settings is future work.

\bibliographystyle{plainnat}
\begin{thebibliography}{10}

\bibitem{lessismore2025}
Less Is More Until It Breaks: Adversarial Vulnerability of VLA Token Compression.
\textit{arXiv:2601.12042}, 2025.

\bibitem{cage2026}
CAGE: Certified Adversarial Robustness Guarantees for Token Compression.
\textit{Proceedings of ICML}, 2026.

\end{thebibliography}

\end{document}
